\documentclass[12pt]{article}
\usepackage[T1]{fontenc}
\usepackage[utf8]{inputenc}
\usepackage{lmodern}
\usepackage{textcomp}
\usepackage{enumitem}
\usepackage{geometry}
\geometry{margin=1in}
\begin{document}
\begin{center}
Answer Key - Economics - Undergraduate Level Assessment 24 \
\small (Answers are ordered exactly as in the question paper.)
\end{center}

\section*{Multiple Choice}
\begin{enumerate}[label=\arabic*.]
\item \textbf{Answer:} A
\\ \textit{Options:} A) Deadweight loss exists; B) Perfect competition; C) Price taker; D) Many sellers and buyers; E) High barriers to entry
\\ \textit{Note:} Deadweight loss exists in a monopoly market because the monopolist restricts output to maximize profits, resulting in a loss of consumer and producer surplus that would occur in a competitive market, whereas monopolistic competition, with its many firms, typically leads to more efficient outcomes with minimal deadweight loss.
\item \textbf{Answer:} E
\\ \textit{Options:} A) A decrease in the cost of tuba-making materials; B) An increase in the number of graduates at tuba maker training school; C) An increase in the tax on tubas; D) An increase in the number of tuba maker training schools; E) An increase in the price of tubas
\\ \textit{Note:} An increase in the price of tubas raises the revenue for tuba makers, which can lead to higher wages as firms seek to attract and retain skilled labor in response to increased demand for their product.
\item \textbf{Answer:} A
\\ \textit{Options:} A) Differences in taxation policies between economies; B) The inclusion or exclusion of foreign direct investment; C) absence of investment; D) Variances in the level of technological advancement; E) Variations in the savings rate of the population
\\ \textit{Note:} Differences in taxation policies between economies can significantly influence consumer spending, investment, and overall economic activity, leading to variations in aggregate demand between nontrading and trading economies.
\item \textbf{Answer:} A
\\ \textit{Options:} A) It introduced mandatory profit sharing between employers and employees; B) It obligated employers to bargain in good faith with a union; C) It outlawed company unions; D) It created a universal pension system for all workers regardless of employer contributions; E) It allowed for automatic union membership for all employees in a workplace
\\ \textit{Note:} The Wagner Act is considered "labor's Magna Charta" because it established fundamental rights for workers to organize and engage in collective bargaining, thus significantly enhancing their economic power and influence in the workplace, akin to the protections and rights granted by the historical Magna Carta.
\item \textbf{Answer:} A
\\ \textit{Options:} A) Personal investments in real estate; B) Buying of stocks and bonds; C) Depreciation of capital; D) Transfer of ownership of existing assets; E) Government spending on infrastructure projects
\\ \textit{Note:} Gross private domestic investment refers to the total amount of investment in physical assets, including personal investments in real estate, which is a significant component of the overall investment in the economy.
\end{enumerate}

\section*{True / False}
\begin{enumerate}[label=\arabic*.]
\setcounter{enumi}{5}
\item \textbf{Answer:} True
\\ \textit{Options:} A) True; B) False
\\ \textit{Note:} Increasing government spending can boost aggregate demand and stimulate economic activity, leading to higher real GDP, especially in a context of underutilized resources, while the absence of immediate tax increases helps to mitigate inflationary pressures in the short run.
\item \textbf{Answer:} False
\\ \textit{Options:} A) True; B) False
\\ \textit{Note:} Lower interest rates in the United States typically lead to a depreciated dollar, which makes U.S. exports cheaper and more competitive abroad, generally resulting in an increase in U.S. exports rather than a decrease.
\item \textbf{Answer:} True
\\ \textit{Options:} A) True; B) False
\\ \textit{Note:} Microeconomics analyzes the behavior and decision-making processes of individual consumers and firms to understand how they interact in markets, which is why its primary focus is on these specific units within the economy.
\item \textbf{Answer:} False
\\ \textit{Options:} A) True; B) False
\\ \textit{Note:} A celebrity's income from endorsements is considered an example of competitive earnings because it reflects their unique skills and marketability, rather than just a transfer of income based on their existing wealth or resources.
\item \textbf{Answer:} False
\\ \textit{Options:} A) True; B) False
\\ \textit{Note:} Economic profits are determined by subtracting both explicit and implicit costs from total revenue, not just subtracting implicit costs from explicit costs.
\end{enumerate}

\section*{Long Form Response}
\begin{enumerate}[label=\arabic*.]
\setcounter{enumi}{10}
\item \textbf{Answer:} The concept of opportunity cost is crucial in understanding Nancy's decision to study for an exam instead of working. In this scenario, Nancy faces an opportunity cost of \$6, which represents the potential earnings she forgoes by not choosing to work. However, the benefits of studying, such as better preparation for the exam and the potential for higher grades, can outweigh this cost. Thus, her choice reflects a careful consideration of the trade-offs between immediate financial gain and long-term academic success. Ultimately, Nancy's decision illustrates how opportunity costs play a vital role in economic decision- making by highlighting the benefits and costs associated with her choices.
\\ \textit{Note:} The answer correctly identifies opportunity cost as the \$6 in potential earnings Nancy gives up to study, while also recognizing that the benefits of improved exam preparation and potential higher grades can justify her decision, illustrating the trade-offs she evaluated.
\item \textbf{Answer:} The Dickey-Fuller test is significant in identifying autocorrelation in time series data, particularly for detecting unit roots that indicate non-stationarity. By assessing whether a time series is stationary, the test helps economists understand the persistence of shocks over time, which is crucial for accurate modeling and forecasting. Identifying autocorrelation up to the third order allows analysts to recognize patterns and dependencies in data, informing decisions related to economic policy and investment strategies. Ultimately, the Dickey-Fuller test enhances the reliability of economic analyses by ensuring that the underlying assumptions of time series models are valid, thereby improving the robustness of conclusions drawn from economic data.
\\ \textit{Note:} correct because the Dickey-Fuller test is essential for detecting non-stationarity in time series data, which allows economists to understand shock persistence and autocorrelation patterns that are vital for effective economic modeling and policy formulation.
\end{enumerate}

\vfill
\noindent\textit{End of Answer Key}
\end{document}